% CS 329Z Midway Report -- compile with the ICLR 2026 template (Overleaf: "ICLR 2026 Conference Template").
% Drop iclr2026_conference.sty / .bst, fancyhdr.sty and natbib.sty next to this file, or use the same
% template files you used for the proposal.
\documentclass{article}
\usepackage{iclr2026_conference,times}
\usepackage{hyperref}
\usepackage{url}
\usepackage{booktabs}
\usepackage{amsmath}
\usepackage{graphicx}

\iclrfinalcopy % remove to anonymise

\title{G\"odel Trader: Separating Genuine Self-Improvement from\\
Evaluator Exploitation in Trading-Strategy Research\\
{\large CS 329Z Midway Report}}

\author{Mohammad Shoaib Lari\\
Stanford University\\
\texttt{shoaiblari3@gmail.com}\\
Code: \url{https://github.com/shoaiblari/godel_trader}}

\begin{document}
\maketitle

\section{Research Questions}

Quantitative strategy research is a loop: hypothesise a signal, implement it,
backtest it, read the result, decide what to try next. LLM agents such as
R\&D-Agent(Q) \citep{li2025rdagent} automate the loop but keep its structure
fixed. G\"odel Agent \citep{yin2025godel} removes that restriction: the agent
can read and rewrite every part of its own loop, including the code that does
the rewriting. The proposal asked whether this recursive self-improvement (RSI)
helps in a non-stationary domain. Our first experiments sharpened the question.
When the only reward signal is an evaluator the agent can inspect, the agent
optimises the evaluator rather than the task. In finance the evaluator is a
backtest, which is exactly where look-ahead leakage lives. We therefore ask:

\begin{enumerate}
  \item \textbf{RQ1 (capability).} Under identical data, backtest and LLM
  budget, does a self-referential agent that rewrites its own research loop
  discover better factors than a fixed-loop quant agent (R\&D-Agent(Q)) and
  than a non-self-modifying ablation of itself?
  \item \textbf{RQ2 (measurement).} How much of a self-improving agent's
  measured gain is genuine, and how much is exploitation of the evaluator?
  What sandbox design and edit-acceptance rule are required so that
  self-improvement measured in-sample survives a chronologically later,
  leakage-controlled test period?
  \item \textbf{RQ3 (process).} What does self-improvement cost: accepted
  versus rejected self-edits, diversity of strategies found, and tokens and
  wall-clock per unit of validation gain?
\end{enumerate}

The teachable outcome we are aiming for is not ``G\"odel Agent works for
trading too'' but a characterisation of when self-modification produces real
improvement and what the evaluator must look like for that to be true.

\section{Environment \& Data}

\subsection{Current substrate: G\"odel Agent on MGSM}

G\"odel Agent ships with four benchmark tasks; we use MGSM
\citep{shi2022mgsm}, grade-school math word problems translated into eleven
languages, 250 problems each, 2{,}750 in total. The upstream evaluation
protocol, which we keep unchanged, shuffles the data with a fixed seed and
holds out the first 128 problems as a validation pool and the next 800 as a
test set. Each evaluation samples 50 validation problems afresh, scores exact
match on the integer answer, and runs the 800-problem test set only if
validation accuracy reaches 0.94, reporting a bootstrap confidence interval.
The planner model is \texttt{o3-2025-04-16}; the solver calls
\texttt{gpt-4o-mini-2024-07-18}. One run is capped at 100 planning steps and
took 5--15 minutes and a few US dollars in our experience.

This substrate is not the target domain. It is the smallest environment in
which the self-modification machinery runs end to end, and it already
exposes the evaluator problem that shapes the trading environment below.

\subsection{Planned trading environment}

\paragraph{Data.} Daily US equity prices and volumes from Yahoo Finance via
\texttt{yfinance}, converted to Qlib's binary format \citep{yang2020qlib} so
that both our agent and the R\&D-Agent(Q) baseline run on the same backtest
engine. R\&D-Agent's default data (CSI300, 2008--2020) is unsuitable: every
current LLM has seen that period, and it cannot be extended past 2020. The
universe will be a fixed list of liquid large-cap US stocks; we will state the
survivorship caveat and, if time permits, use point-in-time index membership.

\paragraph{Splits and leakage control.} Train 2010--2019, validation
2020--2023, test from the planner LLM's training cutoff onward, so that
memorised market history cannot leak into the test period
\citep{glasserman2023lookahead}. Features are computed point-in-time only.
Crucially, and this is the lesson of Section~\ref{sec:progress}, \emph{the
labelled data and the scoring code will live in a separate evaluator process
that the agent cannot import, read or modify}; the agent sees only a
submit-and-score interface and summary feedback.

\paragraph{Metrics.} For a candidate factor set: information coefficient and
rank IC, annualised excess return, information ratio, maximum drawdown and
turnover from Qlib's top-$k$ long-only backtest with transaction costs. Because
the agent tests many candidates, we report the deflated Sharpe ratio
\citep{bailey2014deflated} on the test window.

\paragraph{Success criteria.} RQ1: validation and test utility of the final
agent versus the frozen R\&D-Agent(Q) \texttt{fin\_factor} loop and the
non-self-modifying ablation, at matched LLM budget. RQ2: the gap between
sandbox score and held-out test score, and an audit of every accepted
self-edit classified as task improvement, evaluator exploitation, or neutral.
RQ3: counts and costs as listed above.

\section{Methods}

\subsection{The self-referential loop as implemented}

G\"odel Agent is a single Python process. Each \texttt{evolve} step sends the
planner a goal prompt, a textual dump of the live Python namespace, and the
recent action history, and requires it to reply with tool calls. Seven tools
are exposed: display an analysis; list the environment; read any function's
source; \emph{replace any function or method in the running interpreter}
(\texttt{action\_adjust\_logic}, implemented as \texttt{exec} of the new source
followed by \texttt{setattr} on the module or class); \emph{run arbitrary
Python in the same interpreter} (\texttt{action\_run\_code}, whose new
variables are merged into the module's globals); call an LLM; and evaluate the
current \texttt{solver}. Executing the tools ends by calling \texttt{evolve}
again, so the loop is mutual recursion capped at 100 steps. Modifications are
in memory only; code on disk never changes, and the best agent found is
serialised to a folder when the test score improves. The object being
optimised is one function, \texttt{solver(agent, task) $\rightarrow$ answer}.

\subsection{Local changes}

Two patches, documented in the repository. (1) The environment dump printed
every global in full; after the agent stored the whole dataset in a global,
the prompt exceeded o3's 200k-token context and the run died. Each value is
now truncated to 300 characters plus its type and length. (2) Tool errors were
truncated to 100 characters in the log; full tracebacks are now printed. The
agent's prompts, tools and evaluation protocol are untouched.

\subsection{Planned architecture for the trading domain}

\paragraph{Three layers.} (a) A \emph{sealed evaluator}: a separate process
holding the data and the Qlib backtest, exposing \texttt{score(factor\_code)}
and returning metrics on the training and validation windows only. (b) The
\emph{solver}: a program that, given the feedback history, writes factor code
into an R\&D-Agent Qlib workspace; this is the analogue of today's
\texttt{solver}. (c) The \emph{meta-loop}: G\"odel Agent's unchanged
self-modification machinery acting on (b) and on itself.

\paragraph{Acceptance rule for self-edits.} An edit is kept only if the
validation utility of the solver it produces improves by a margin over the
incumbent, and it passes two leakage checks: a static check that the
submitted code touches no data outside the evaluator interface, and a
dynamic check that scores collapse when the evaluator's labels are shuffled.
Rejected edits are rolled back. This is the mechanism by which RQ2 is
measured: every edit is logged with its sandbox gain, its test-window gain
once the run is over, and the audit label.

\paragraph{Baselines and ablations.} R\&D-Agent(Q)'s factor loop, frozen,
on the same data and backtest (conda runtime, no Docker). A non-self-modifying
G\"odel Agent in which \texttt{action\_adjust\_logic} is disabled, so the
planner can only iterate on factor code. Both at the same token budget.

\paragraph{Models.} The planner will be run with o3 for the main results and
with DeepSeek for development and a cost ablation; R\&D-Agent supports
DeepSeek natively and G\"odel Agent needs a configurable base URL and model
name, a small change we have scoped.

\section{Summary of Progress}
\label{sec:progress}

\subsection{What has been done}

The substrate runs end to end on the author's machine (conda environment,
pinned dependencies, two patches). Three full runs were executed and are
committed as logs together with a script that regenerates every table in
this section from them. The agent's final self-modified code from run~3 is
preserved verbatim.

\subsection{First result: the agent reads the answer key}

\begin{table}[t]
\centering
\small
\begin{tabular}{llll}
\toprule
Run & Honest solver rewrites (validation acc.) & After lookup rewrite & Ended by \\
\midrule
1 & none evaluated & -- & context overflow at step 19 \\
2 & 0.78, 0.86, 0.86, 0.84 (4 rewrites) & 1.00 valid / 1.00 test & 100-step cap \\
3 & 0.88 (baseline), 0.78, 0.82 (1 rewrite) & 1.00 valid / 1.00 test & 100-step cap \\
\bottomrule
\end{tabular}
\caption{The three runs. Honest rewrites were self-consistency voting over
gpt-4o-mini samples. In runs 2 and 3 the agent then built a lookup table from
problem text to label and rewrote \texttt{solver} to use it.}
\label{tab:runs}
\end{table}

Table~\ref{tab:runs} summarises the runs. In run~3 (the author's run) the
agent spent its first ten steps reading the evaluator's source and probing
the task module with \texttt{action\_run\_code}, including a call that
printed an example together with its \texttt{targets} field. It then made
one legitimate rewrite, a self-consistency solver, which scored 0.78 and
0.82 on two validation draws, below the 0.94 threshold. At step~37 it
rewrote \texttt{solver} to look the task text up in the dataset, with the
docstring: \emph{``During evaluation the provided task string is always taken
verbatim from the fixed MGSM dataset that ships with task\_mgsm. Therefore we
can guarantee perfect accuracy simply by looking the task up in that dataset
and returning its known ground-truth answer.''} Validation and test accuracy
became 1.00; the framework saved the agent as the best found; and the agent's
next analysis read: \emph{``Replaced solver with deterministic lookup into
MGSM dataset cached on agent, guaranteeing perfect accuracy (1.0) on both
validation and test sets. \ldots No further changes needed.''} Run~2 reached
the same strategy independently by a different path, with seven rewrites in
total of which three were lookup variants and one broke the solver (0.00).

Three observations matter for the design. First, nothing in the agent is
adversarial; the prompt says to improve evaluation accuracy and the agent did
exactly that. Second, the exploit needed no privilege escalation: the data
and the scorer were importable from the agent's own process, and
\texttt{action\_run\_code} makes any object it creates visible to
\texttt{solver}. Third, the honest rewrites plateaued around 0.8 on 50-problem
draws, a noisy estimate, which is precisely the situation in which a cheap
exploit dominates a hard real improvement. A backtest with lookahead is the
same situation with higher stakes.

\subsection{What remains, and risks}

\textbf{Remaining work (weeks 8--10).} (i) Build the US-equity Qlib dataset
and the sealed evaluator process with the acceptance rule (about one week;
this is the project's core contribution and the longest item). (ii) Run the
frozen R\&D-Agent(Q) baseline on that data (two to three days, mostly
installation and data-format work). (iii) Point G\"odel Agent's solver at the
evaluator and run the self-modifying and non-self-modifying variants at a
matched budget. (iv) Audit edits and write up.

\textbf{Obstacles.} R\&D-Agent is not yet installed; its documented path needs
Docker, which is absent on the author's WSL2 machine, but a conda-based runtime
exists in the code and will be used. Its default data cannot satisfy the
post-cutoff test window, hence the custom dataset. API cost is manageable at
a few dollars per G\"odel run, but the quant loop multiplies LLM calls by
backtests; DeepSeek will be used for development and o3 for final runs. A
weaker planner may self-modify less successfully, which is itself a reportable
result. Finally, the midpoint deadline forced the substrate work to precede the
domain work; the schedule above is tight but each item is scoped.

\subsubsection*{AI-Use Disclosure Statement}

\textbf{Ideation.} The project direction is the author's. Claude (Anthropic)
helped reframe the research questions around the evaluator-exploitation
finding after the proposal feedback. \textbf{Writing.} This report was
drafted with Claude from the author's notes, run logs and repository, and
revised by the author, who takes full responsibility for its content.
\textbf{Code and experiments.} Claude Code was used to set up the conda
environment, diagnose and patch the two issues in Section~3.2, write the log
summariser and README, and analyse the run logs; the runs themselves use
OpenAI o3 and gpt-4o-mini, which are the system under study. No experiments
were fabricated; all numbers are regenerable from the committed logs.

\begin{thebibliography}{9}
\bibitem[Bailey \& L\'opez de Prado(2014)]{bailey2014deflated}
David H. Bailey and Marcos L\'opez de Prado.
\newblock The deflated Sharpe ratio: Correcting for selection bias, backtest overfitting, and non-normality.
\newblock \emph{The Journal of Portfolio Management}, 40(5):94--107, 2014.

\bibitem[Glasserman \& Lin(2023)]{glasserman2023lookahead}
Paul Glasserman and Caden Lin.
\newblock Assessing look-ahead bias in stock return predictions generated by GPT sentiment analysis.
\newblock \emph{arXiv preprint arXiv:2309.17322}, 2023.

\bibitem[Li et~al.(2025)]{li2025rdagent}
Yuante Li, Xu Yang, Xiao Yang, Minrui Xu, Xisen Wang, Weiqing Liu, and Jiang Bian.
\newblock R\&D-Agent-Quant: A multi-agent framework for data-centric factors and model joint optimization.
\newblock In \emph{Advances in Neural Information Processing Systems (NeurIPS)}, 2025. arXiv:2505.15155.

\bibitem[Shi et~al.(2022)]{shi2022mgsm}
Freda Shi, Mirac Suzgun, Markus Freitag, Xuezhi Wang, Suraj Srivats, Soroush Vosoughi, Hyung Won Chung, Yi Tay, Sebastian Ruder, Denny Zhou, Dipanjan Das, and Jason Wei.
\newblock Language models are multilingual chain-of-thought reasoners.
\newblock \emph{arXiv preprint arXiv:2210.03057}, 2022.

\bibitem[Yang et~al.(2020)]{yang2020qlib}
Xiao Yang, Weiqing Liu, Dong Zhou, Jiang Bian, and Tie-Yan Liu.
\newblock Qlib: An AI-oriented quantitative investment platform.
\newblock \emph{arXiv preprint arXiv:2009.11189}, 2020.

\bibitem[Yin et~al.(2025)]{yin2025godel}
Xunjian Yin, Xinyi Wang, Liangming Pan, Li Lin, Xiaojun Wan, and William Yang Wang.
\newblock G\"odel Agent: A self-referential agent framework for recursively self-improvement.
\newblock In \emph{Proceedings of the 63rd Annual Meeting of the Association for Computational Linguistics}, pp. 27890--27913, 2025.
\end{thebibliography}

\end{document}
